\chapter{Trabalhos Relacionados}
\section{Aprendizado de Máquina com Emissão Acústica para Detecção de Falhas}

\subsection{Introdução}

A Emissão Acústica (EA) é uma técnica consolidada para a detecção de falhas em materiais e estruturas. Nos últimos anos, o uso de aprendizado de máquina (ML, do inglês Machine Learning) em combinação com EA tem se mostrado promissor, especialmente para melhorar a precisão e a automação da detecção de falhas. O aprendizado de máquina permite identificar padrões complexos em dados acústicos que seriam difíceis de detectar manualmente, tornando-o uma abordagem inovadora para o monitoramento preditivo e a manutenção de ativos críticos.

\subsection{História do Uso de Aprendizado de Máquina com Emissão Acústica}

A combinação de aprendizado de máquina e EA para detecção de falhas é relativamente recente. O uso de EA começou a se popularizar na década de 1950, mas foi apenas com os avanços em processamento de sinais e inteligência artificial, a partir das décadas de 1990 e 2000, que surgiram as primeiras tentativas de integrar técnicas de aprendizado de máquina ao processo de análise de sinais de EA (Pollock, 2012).
O desenvolvimento de algoritmos supervisionados, como redes neurais artificiais, e de técnicas não supervisionadas, como agrupamento (clustering), possibilitou a criação de sistemas automáticos de detecção de falhas baseados em EA. Mais recentemente, com o aumento da capacidade computacional e o advento de métodos mais avançados de aprendizado, como deep learning, a análise de grandes volumes de dados acústicos tornou-se mais eficiente e precisa (Sohn et al., 2010).

\subsection{Fundamentos do Aprendizado de Máquina e Emissão Acústica}

A Emissão Acústica é um método não destrutivo que detecta ondas acústicas geradas por eventos internos em materiais. Essas ondas, ao serem captadas por sensores piezoelétricos, geram sinais complexos, cuja análise pode fornecer informações sobre a presença e a localização de falhas, como trincas, corrosão ou fadiga.

Com o aprendizado de máquina, esses sinais de EA podem ser processados de maneira mais robusta. As técnicas de ML são usadas para:


Classificação: Identificar automaticamente se os eventos acústicos indicam uma falha ou se são ruídos ou outros eventos normais.


Detecção de Anomalias: Sistemas não supervisionados de aprendizado de máquina podem ser treinados para detectar padrões anômalos em dados de EA, indicando a possível presença de falhas incipientes.


Previsão: Técnicas de regressão e redes neurais são aplicadas para prever a progressão de uma falha com base no histórico de sinais de EA.
Essas abordagens têm uma vantagem significativa em relação aos métodos tradicionais de análise de sinais, pois podem aprender a identificar padrões complexos sem a necessidade de definir previamente regras rígidas, o que as torna mais adaptáveis a diferentes materiais e condições operacionais (Aggelis et al., 2012).

\subsection{Aplicações de Aprendizado de Máquina com Emissão Acústica}

Monitoramento Estrutural em Tempo Real
Um dos principais campos de aplicação é o monitoramento de grandes estruturas, como pontes e edifícios. Sensores de EA acoplados a sistemas de aprendizado de máquina podem detectar a propagação de trincas em tempo real. A capacidade de ML de reconhecer padrões complexos em sinais de EA ajuda a distinguir eventos críticos de ruídos de fundo, fornecendo alertas precoces para intervenções (Sharma et al., 2021).

\subsubsection{Indústria Aeroespacial}

Na indústria aeroespacial, o aprendizado de máquina tem sido usado para analisar sinais de EA captados de componentes críticos, como fuselagens e asas de aeronaves. Modelos de ML treinados com dados históricos de falhas podem identificar possíveis defeitos causados por fadiga de materiais, garantindo maior segurança (Prosser, 1991).

\subsubsection{Manutenção Preditiva de Máquinas Rotativas}

Máquinas rotativas, como turbinas, motores e bombas, também se beneficiam dessa integração. O aprendizado de máquina pode classificar sinais de EA captados de rolamentos e eixos para identificar desgaste, desalinhamento e cavitação. Ao usar técnicas de detecção de anomalias, os sistemas podem prever falhas antes que elas ocorram, reduzindo o tempo de inatividade e custos de manutenção (Gupta et al., 2017).

\subsubsection{Monitoramento de Corrosão}

Sistemas de ML aplicados à detecção de corrosão usando EA são capazes de identificar o início de processos corrosivos em materiais metálicos, como em tubulações e estruturas offshore. A sensibilidade desses sistemas permite detectar os primeiros sinais de corrosão por pite, ajudando na preservação de infraestruturas críticas (Chaudhry et al., 1997).

\subsection{Desafios e Oportunidades}

Apesar dos avanços, a aplicação de aprendizado de máquina em EA enfrenta alguns desafios. Um dos principais problemas é a necessidade de grandes volumes de dados etiquetados para treinar modelos supervisionados. A aquisição de tais dados pode ser cara e demorada, especialmente em ambientes industriais onde as falhas ocorrem com pouca frequência. Outro desafio é a presença de ruídos, que pode comprometer a precisão dos modelos de ML.
Com o aumento da capacidade de coleta e processamento de dados em tempo real, há uma oportunidade crescente de aplicar modelos mais complexos, como redes neurais profundas, para melhorar a detecção de falhas. Além disso, o uso de técnicas de aprendizado por reforço e aprendizado federado pode permitir a criação de sistemas mais adaptáveis e distribuídos, que aprendem continuamente à medida que novos dados são coletados (Zhu et al., 2020).

\subsubsection{Conclusão}
A combinação de aprendizado de máquina e Emissão Acústica representa uma evolução significativa no campo da detecção de falhas. Essa integração permite monitoramento em tempo real, predição precisa e automação de processos, tornando-a uma ferramenta valiosa para a indústria 4.0 e para sistemas de manutenção preditiva. Com o avanço contínuo dos algoritmos de ML e das tecnologias de sensoriamento, espera-se que essa abordagem se torne ainda mais robusta e amplamente aplicável em uma variedade de indústrias.

% 3.8.	Referências
% Aggelis, D. G., Verbruggen, S., Tsangouri, E., & Tysmans, T. (2012). Monitoring of crack propagation in a concrete beam during four-point bending by combined techniques. NDT & E International, 48, 66-72. 
% Chaudhry, Z., Gan, D. D., & Raghavan, V. (1997). Acoustic emission monitoring of corrosion in metallic structures. Journal of Engineering Materials and Technology, 119(2), 192-198. 
% Gupta, M. S., Behera, D. K., & Nayak, A. K. (2017). Acoustic emission-based condition monitoring of rotating machines: A review. Engineering Failure Analysis, 79, 626-642.

% Pollock, A. A. (2012). Fundamentals of acoustic emission testing. Journal of Acoustic Emission, 30(2), 55-69. 
% Prosser, W. H. (1991). Advanced AE techniques in composite materials research. Journal of Acoustic Emission, 9(4), 292-306.
% Sharma, N., Chakraborty, S., & Chowdhury, R. (2021). Fault detection using machine learning algorithms in structural health monitoring: A review. Journal of Civil Structural Health Monitoring, 11, 1065-1091.
% Sohn, H., Dutta, D., & Taylor, S. G. (2010). Acoustic emission monitoring for crack detection and localization in plate structures. Nondestructive Testing and Evaluation, 25(4), 321-331.
% Zhu, M., Wang, Z., Xu, Y., & Zhang, T. (2020). An overview of federated learning. Journal of Nondestructive Testing, 40, 1231-1245.
